\section{Shapes and structures}\label{sec:shapes}

\subsection{Small categories as shapes}

Recall that a \emph{category} $C$ has objects, and for any two objects $U,V$ a set $\Hom(V,U)$ of arrows $f : V\to U$, with an identity $\id_U : U\to U$ and an associative composition $g\circ f$ for $f : W\to V$, $g : V\to U$. The category is \emph{small} when its objects form a set. We shall think of a small category as a \emph{shape}: a pattern of \q{places} (the objects) and of ways of \q{looking from one place at another} (the arrows).

\begin{example}[A monoid]\label{ex:BM}
Let $M$ be a monoid, with unit $1$. The category $\B M$ has a single object $\ast$; its arrows $\ast\to\ast$ are the elements of $M$; composition is multiplication, $g\circ f = gf$, and $\id_\ast = 1$. When $M = G$ is a group, every arrow of $\B G$ is invertible.
\end{example}

\begin{example}[A preorder]\label{ex:preorder}
A \emph{preorder} on a set $P$ is a reflexive and transitive relation $\le$. It gives a category whose objects are the elements of $P$, with exactly one arrow $p\to q$ when $p\le q$ and none otherwise. Composition exists by transitivity, identities by reflexivity. Examples: $(\mathbb R,\le)$; the subsets of a set ordered by inclusion; the positive integers ordered by divisibility; the chain $\omega = \{0<1<2<\cdots\}$; the \q{V} shape $\{a<b,\ a<c\}$.
\end{example}

\begin{example}[The shape of graphs]\label{ex:graphshape}
The category $\mathsf{G}$ has two objects $V$ (\q{vertex}) and $E$ (\q{edge}), and besides the identities two arrows $s,t : V\to E$ (\q{source}, \q{target}). Picture: $V\rightrightarrows E$.
\end{example}

\begin{example}[The simplex category]\label{ex:Delta}
$\Delta$ has as objects the finite ordinals $[n] = \{0<1<\cdots<n\}$ for $n\ge 0$, and as arrows the non-decreasing maps. The \emph{face maps} $\delta^i : [n-1]\to[n]$ skip the value $i$; for instance $\delta^0,\delta^1 : [0]\to[1]$ pick the vertices $1$ and $0$ of the segment $[1]$. The \emph{augmented} simplex category $\Delta_+$ adds the empty ordinal $[-1] = \emptyset$.
\end{example}

\begin{example}[Finite sets]\label{ex:finsets}
$\mathsf{Fin}$: finite sets and all maps. $\mathsf{FI}$: finite sets and injective maps. $\mathsf{Fin}_{>0}$: non-empty finite sets and all maps.
\end{example}

A category is a \emph{groupoid} if every arrow is invertible: $\B G$ for a group $G$, a preorder in which $\le$ is symmetric (an equivalence relation), the category of finite sets and bijections.

\subsection{Presheaves: structures of a given shape}

\begin{definition}[Presheaf]
A \emph{presheaf} $F$ on a small category $C$ gives, for each object $U$, a set $F(U)$ (the \emph{elements of $F$ at $U$}), and for each arrow $f : V\to U$ a map $F(U)\to F(V)$, written $x\mapsto\res xf$ and called \emph{restriction along $f$}, such that
\[
\res x{\id_U} = x,\qquad \res{(\res xg)}{f} = \res x{g\circ f}\quad(f : W\to V,\ g : V\to U).
\]
A \emph{morphism} of presheaves $\varphi : F\to F'$ is a family of maps $\varphi_U : F(U)\to F'(U)$ commuting with restrictions: $\varphi_V(\res xf) = \res{\varphi_U(x)}{f}$. We write $\PSh(C)$ for the category of presheaves on $C$.
\end{definition}

In the language of categories, a presheaf is a functor $C^{\op}\to\Set$: arrows go backwards, because restriction goes from $U$ to $V$ when $f$ goes from $V$ to $U$. The word comes from topology, where $F(U)$ is a set of functions on an open set $U$ and restriction is restriction of functions to a smaller open set. The point of this text is that many familiar structures are presheaves.

\begin{example}[Right $M$-sets]\label{ex:Msets}
A presheaf on $\B M$ is a set $X = F(\ast)$ with, for each $m\in M$, a map $x\mapsto\res xm$; writing $x\cdot m$ for $\res xm$, the two axioms read $x\cdot 1 = x$ and $(x\cdot g)\cdot f = x\cdot(gf)$. So a presheaf on $\B M$ is exactly a \emph{right $M$-set}, a set with a right action of $M$, and a morphism is an equivariant map. For a group $G$, these are the $G$-sets of group theory.
\end{example}

\begin{example}[Automata]\label{ex:automata}
Let $\Sigma$ be a finite alphabet and $\Sigma^*$ the free monoid of words on $\Sigma$, with concatenation. A right $\Sigma^*$-set is a set $X$ with a map $X\times\Sigma\to X$, $(x,a)\mapsto x\cdot a$, extended to words by $x\cdot(a_1a_2\cdots a_k) = (\cdots((x\cdot a_1)\cdot a_2)\cdots)\cdot a_k$. This is a \emph{deterministic automaton} without initial or final states: $X$ is the set of states, and reading the letter $a$ in state $x$ moves to state $x\cdot a$. For a one-letter alphabet, $\Sigma^*\cong(\mathbb N,+)$, and an automaton is a set with a single map $\tau : X\to X$: a \emph{discrete dynamical system}.
\end{example}

\begin{example}[Directed graphs]\label{ex:graphs}
A presheaf $F$ on the shape $\mathsf G$ of Example~\ref{ex:graphshape} is a set of vertices $F(V)$, a set of edges $F(E)$, and two maps $\res{-}{s},\res{-}{t} : F(E)\to F(V)$ giving the source and the target of each edge. This is a directed graph (multiple edges and loops allowed).
\end{example}

\begin{example}[Simplicial sets]
A presheaf on $\Delta$ is a \emph{simplicial set}: a set $X_n$ of $n$-simplices for each $n$, with face maps $X_n\to X_{n-1}$ (restriction along $\delta^i$) and degeneracy maps. Simplicial sets are the combinatorial models of spaces used in algebraic topology. Presheaves on $\Delta_+$ are augmented simplicial sets; presheaves on $\mathsf{Fin}_{>0}$ are sometimes called symmetric simplicial sets.
\end{example}

\begin{example}[Families indexed by a preorder]\label{ex:preorderpsh}
Let $P$ be a preorder. A presheaf on $P^{\op}$ (the opposite preorder) gives a set $F(p)$ for each $p$ and a map $F(p)\to F(q)$ whenever $p\le q$, compatible with composition: a \emph{family of sets indexed by $P$, varying covariantly}. Equivalently, a functor $P\to\Set$. For $P = \omega^{\op}$, a presheaf on $\omega$ is a sequence of sets and maps $X_0\leftarrow X_1\leftarrow X_2\leftarrow\cdots$; this is the \emph{topos of trees} of computer science \cite{birkedal}, where $X_n$ is thought of as \q{what is known after $n$ steps}.
\end{example}

\begin{exercise}\label{exo:reflgraph}
A \emph{reflexive graph} is a directed graph with a chosen loop at each vertex. Find a category $\mathsf G_r$ with two objects such that reflexive graphs are presheaves on $\mathsf G_r$. (Hint: add an arrow $E\to V$.)
\end{exercise}

\subsection{Representable presheaves and sieves}

Every object $U$ of $C$ gives a presheaf $\y U$, the presheaf \emph{represented} by $U$:
\[
(\y U)(V) = \Hom(V,U),\qquad \res{g}{f} = g\circ f.
\]
In words, the elements of $\y U$ at $V$ are the ways of \q{looking from $V$ at $U$}, and restricting is precomposing. For $\B M$, $\y\ast$ is $M$ acting on itself by right multiplication. For the graph shape, $\y E$ is the graph with one edge (and its two endpoints), and $\y V$ the graph with one vertex.

\begin{lemma}[Yoneda, the easy half]\label{lem:yoneda}
For a presheaf $F$ and an object $U$, a morphism $\varphi : \y U\to F$ is determined by the single element $x = \varphi_U(\id_U)\in F(U)$, through $\varphi_V(f) = \res xf$; and every $x\in F(U)$ arises in this way.
\end{lemma}

\begin{proof}
If $\varphi$ is a morphism, $\varphi_V(f) = \varphi_V(\res{\id_U}f) = \res{\varphi_U(\id_U)}f$. Conversely, $f\mapsto\res xf$ commutes with restrictions by the axiom $\res{(\res xg)}f = \res x{g\circ f}$.
\end{proof}

A sub-presheaf of $\y U$ is called a \emph{sieve} on $U$: it is a set $S$ of arrows into $U$ such that $f\in S$ implies $f\circ h\in S$ for every $h$. The sieve \emph{generated} by an arrow $f : V\to U$ is the set of arrows that factor through $f$, that is, of the form $f\circ h$.

\subsection{Why presheaves form a topos}

The category $\PSh(C)$ has all the constructions one is used to in $\Set$: products, disjoint unions, quotients, function spaces, and a \q{set of truth values} (the subobject classifier). It is the first example of a \emph{Grothendieck topos}; we shall meet more general ones in \S\ref{sec:sites}. We do not need the definition of a topos in what follows: for presheaf topoi, everything can be computed element by element, as we now do for subobjects and negation.
